\documentclass[aspectratio=169,11pt]{beamer}
\usepackage{../../../shared/theme/esmad_beamer_theme}
\usepackage{amsmath,amssymb}
\usepackage{booktabs}

\title[Causal Mediation]{Causal Mediation Analysis}
\subtitle{Direct Effects, Indirect Effects, and Mechanisms}
\date{\today}

\newcommand{\E}{\mathbb{E}}

\begin{document}

\begin{frame}
\titlepage
\end{frame}

\begin{frame}{Learning Goals}
\begin{itemize}
\item distinguish total, direct, and indirect causal effects;
\item define mediator potential outcomes and natural-effect estimands;
\item explain sequential ignorability and mediator-outcome confounding;
\item recognize post-treatment confounding as a major identification threat;
\item distinguish natural effects from interventional effects;
\item use sensitivity analysis when mediator assumptions are uncertain.
\end{itemize}
\end{frame}

\begin{frame}{Outline}
\tableofcontents
\end{frame}

\section{Why Mediation?}

\begin{frame}{From “Does It Work?” to “How Does It Work?”}
Let treatment \(D\) affect outcome \(Y\) partly through mediator \(M\):
\[
D \rightarrow M \rightarrow Y.
\]

There may also be a direct path:
\[
D \rightarrow Y.
\]

\begin{block}{Mechanism question}
How much of the treatment effect operates through \(M\)?
\end{block}
\end{frame}

\begin{frame}{Total Effect}
The total effect is
\[
TE=\E[Y(1)-Y(0)].
\]

It combines all causal pathways from treatment to outcome.
\end{frame}

\section{Mediator Potential Outcomes}

\begin{frame}{Potential Mediator}
Let
\[
M(d)
\]
denote the mediator value that would occur under treatment \(d\).

Let
\[
Y(d,m)
\]
denote the outcome if treatment were set to \(d\) and mediator to \(m\).
\end{frame}

\begin{frame}{Nested Counterfactuals}
Natural effects use nested potential outcomes such as
\[
Y(1,M(0)).
\]

This asks what the outcome would be under treatment, while forcing the mediator to whatever value it would naturally have taken under control.
\end{frame}

\section{Natural Direct and Indirect Effects}

\begin{frame}{Natural Direct Effect}
One version of the natural direct effect is
\[
NDE
=
\E[
Y(1,M(0))
-
Y(0,M(0))
].
\]

The mediator is held at its natural control value while treatment changes.
\end{frame}

\begin{frame}{Natural Indirect Effect}
A natural indirect effect can be written
\[
NIE
=
\E[
Y(1,M(1))
-
Y(1,M(0))
].
\]

Treatment is fixed while the mediator changes from its natural control value to its natural treated value.
\end{frame}

\begin{frame}{Decomposition}
Under the chosen definitions,
\[
TE=NDE+NIE.
\]

\begin{alertblock}{Caution}
This decomposition is causal only under assumptions strong enough to identify the nested counterfactuals.
\end{alertblock}
\end{frame}

\section{Controlled Direct Effects}

\begin{frame}{Controlled Direct Effect}
For a fixed mediator level \(m\),
\[
CDE(m)
=
\E[Y(1,m)-Y(0,m)].
\]

This asks what treatment would do if the mediator could be externally fixed to \(m\).
\end{frame}

\begin{frame}{Why Controlled and Natural Effects Differ}
Controlled effects intervene on a specific mediator value.

Natural effects intervene on whichever mediator value a unit would naturally attain under another treatment condition.

These are different scientific questions.
\end{frame}

\section{Identification Assumptions}

\begin{frame}{Treatment-Outcome Exchangeability}
A standard assumption is that, conditional on pre-treatment covariates \(X\),
\[
Y(d,m) \perp D \mid X.
\]

Randomized treatment can make this assumption plausible.
\end{frame}

\begin{frame}{Mediator-Outcome Exchangeability}
We also need a condition such as
\[
Y(d,m) \perp M \mid D,X.
\]

This is much harder: the mediator is usually not randomized.
\end{frame}

\begin{frame}{Sequential Ignorability}
A common identification framework combines:
\begin{itemize}
\item no unmeasured treatment-outcome confounding;
\item no unmeasured treatment-mediator confounding;
\item no unmeasured mediator-outcome confounding after conditioning on treatment and baseline covariates.
\end{itemize}
\end{frame}

\section{Post-Treatment Confounding}

\begin{frame}{A Serious Failure Mode}
Suppose treatment changes a variable \(L\), and \(L\) affects both mediator and outcome:
\[
D \rightarrow L \rightarrow M,
\]
and
\[
L \rightarrow Y.
\]

Then \(L\) is a treatment-induced mediator-outcome confounder.
\end{frame}

\begin{frame}{Why Ordinary Adjustment Can Fail}
Conditioning on \(L\) can block part of the treatment effect while leaving mediation estimands unidentified.

\begin{alertblock}{Core lesson}
Post-treatment covariates cannot be handled like ordinary baseline confounders.
\end{alertblock}
\end{frame}

\section{Exposure-Mediator Interaction}

\begin{frame}{Treatment Can Change the Mediator Effect}
Suppose
\[
Y
=
\beta_0
+
\beta_1 D
+
\beta_2 M
+
\beta_3 D M
+
\varepsilon.
\]

If \(\beta_3\neq 0\), the mediator effect depends on treatment status.
\end{frame}

\begin{frame}{Implication}
Simple “product of coefficients” mediation formulas generally fail when:
\begin{itemize}
\item treatment interacts with mediator;
\item models are nonlinear;
\item treatment or mediator is binary;
\item confounding structure is complex.
\end{itemize}
\end{frame}

\section{Interventional Effects}

\begin{frame}{Why Interventional Effects?}
Natural effects rely on nested counterfactuals that can be scientifically awkward.

Interventional effects instead define stochastic interventions on the mediator distribution.
\end{frame}

\begin{frame}{Interventional Indirect Effect}
Rather than forcing each unit to its own counterfactual mediator value, one can compare outcomes under mediator draws from treatment-specific mediator distributions.

\begin{block}{Benefit}
This can avoid some cross-world interpretations and may remain meaningful with treatment-induced mediator-outcome confounding under suitable formulations.
\end{block}
\end{frame}

\section{Estimation}

\begin{frame}{Regression-Based Mediation}
A simple workflow models:
\[
M=f(D,X)+\eta
\]
and
\[
Y=g(D,M,X)+\varepsilon.
\]

Causal interpretation depends on identification assumptions, not on the fact that two regressions were fitted.
\end{frame}

\begin{frame}{Simulation / G-Computation}
A more general approach:
\begin{enumerate}
\item model the mediator distribution;
\item model the outcome conditional on treatment, mediator, and covariates;
\item simulate mediator and outcome counterfactuals under the target intervention;
\item average the resulting causal contrasts.
\end{enumerate}
\end{frame}

\section{Sensitivity Analysis}

\begin{frame}{Unmeasured Mediator-Outcome Confounding}
Mediator-outcome confounding is often the most fragile assumption.

Sensitivity analysis asks how strong an omitted common cause would need to be to materially change the indirect effect.
\end{frame}

\begin{frame}{Correlation-Sensitivity Intuition}
In linear settings, one can parameterize correlation between mediator-model and outcome-model disturbances.

As that residual correlation increases, recompute direct and indirect effects.
\end{frame}

\section{Applied Reasoning}

\begin{frame}{Example: Marketing Intervention}
Suppose an email campaign \(D\) changes site engagement \(M\), which then changes purchasing \(Y\).

Potential issues:
\begin{itemize}
\item prior engagement confounds mediator and outcome;
\item treatment changes another behavior that also affects purchasing;
\item mediator measurement is noisy;
\item treatment changes the effect of engagement on purchasing.
\end{itemize}
\end{frame}

\begin{frame}{A Credible Mediation Workflow}
\begin{enumerate}
\item Define the total, direct, and indirect estimands.
\item Draw the causal graph including post-treatment variables.
\item Separate baseline from treatment-induced confounders.
\item State the mediation identification assumptions.
\item Choose natural or interventional effects deliberately.
\item Fit mediator and outcome models.
\item allow treatment-mediator interactions where plausible.
\item run sensitivity analysis for mediator-outcome confounding.
\item avoid interpreting regression decomposition as mechanism proof.
\end{enumerate}
\end{frame}

\begin{frame}{Common Failure Modes}
\begin{itemize}
\item treating association between mediator and outcome as mediation;
\item adjusting mechanically for the mediator and calling the remaining coefficient “direct”;
\item ignoring treatment-induced confounding;
\item using product-of-coefficients formulas outside their assumptions;
\item interpreting natural effects without acknowledging cross-world assumptions;
\item skipping sensitivity analysis for mediator-outcome confounding.
\end{itemize}
\end{frame}

\begin{frame}{Synthesis: Mechanisms Need Stronger Assumptions}
\begin{itemize}
\item The total effect asks whether treatment changes the outcome.
\item Mediation asks how the effect operates through intermediate variables.
\item Natural direct and indirect effects use nested counterfactuals.
\item Controlled direct effects answer a different intervention question.
\item Mediator-outcome confounding is typically not removed by treatment randomization.
\item Treatment-induced confounders complicate or invalidate standard mediation formulas.
\item Interventional effects provide alternative mechanism estimands.
\item Sensitivity analysis is essential when mediation assumptions are fragile.
\end{itemize}
\end{frame}

\begin{frame}{Selected References}
\small
\begin{itemize}
\item Robins, J. M., and Greenland, S. (1992). Identifiability and Exchangeability for Direct and Indirect Effects. \textit{Epidemiology}, 3(2), 143--155.
\item Pearl, J. (2001). Direct and Indirect Effects. \textit{Proceedings of UAI}.
\item Imai, K., Keele, L., and Tingley, D. (2010). A General Approach to Causal Mediation Analysis. \textit{Psychological Methods}, 15(4), 309--334.
\item VanderWeele, T. J. (2015). \textit{Explanation in Causal Inference}. Oxford University Press.
\item Vansteelandt, S., and Daniel, R. M. (2017). Interventional Effects for Mediation Analysis with Multiple Mediators. \textit{Epidemiology}, 28(2), 258--265.
\end{itemize}
\end{frame}

\contactslide

\end{document}
